\newpage
\begin{center}
    \Huge Third Warning to McDowell \\
    \Large (Red Card)
\end{center}


So there you have it. We've come to the point in our narrative where Rex Stout gathers all his characters into Nero Wolfe's office so he can unmask the killer---causing an especially emotional reader to cry, ``How come \emph{I} didn't figure that out?!'' In accordance with the rules of a classical detective story, the author has no aces hidden up his sleeve; the reader has been provided with every last shred of evidence, to the point where I'm sure many of them have already pieced this puzzle together themselves.

As for the rest, their only option is to acquaint themselves with a certain manuscript, whose authenticity admittedly has yet to be confirmed by experts. For almost twenty centuries, it lay hidden in a sealed jug, which speleologists from Ben-Gurion University recently found by chance while exploring one of the karst caves near Jerusalem. The manuscript's author had actually been summoned from oblivion once before by the sheer power of Bulgakov's genius. As it now turns out, he led an independent existence altogether, surprisingly retaining all the qualities the Master once endowed him with. And so, I yield the floor to the ``man who never parted with his hood,''\footnote{\emph{The Master and Margarita}, p. 273} the chief of the secret service for the procurator of Judea, the military tribune Afranius.\footnote{In his paper ``Observations on the Motif Structure of Bulgakov's Master and Margarita'' (Daugava, No. 1, 1989, pp. 78-81), Boris Gasparov appeals to a number of textological arguments to directly identify Afranius with… Woland. In his opinion, ``chief of the secret police'' is nothing more than Satan's mask in Yershalaim, in the same way that ``professor of black magic'' serves as his mask in Moscow. However, there are a number of serious objections one can make to such an identification. I, for one, can imagine Woland---Lermontov's ``free son of the ether''---in a variety of different hypostases, but certainly not as a member of the service class; Afranius, meanwhile, speaks quite definitively on this subject: ``I have worked in Judea for fifteen years, Procurator. I began my service under Valerius Gratus'' (p. 273). From my point of view, Gasparov's version is interesting mainly as a reflection of modern man's particular tendency to demonize the secret services. Andrzej Wajda's postmodern Afranius from Pilate and Others, who has traded out his hood and sword for a pair of shades and a pistol in a shoulder holster, seems far more compelling by comparison. (K.E.)}